\subsection{Finite-state automata}

\begin{definition}
	A finite-state automaton $\mathcal{A}$ is a tuple $\langle \Sigma, Q,
	I, F, \Delta \rangle$ where $\Sigma$ is an alphabet, $Q$ is a finite set
	of states, $I \subseteq Q$ is a set of initial states, $F \subseteq Q$
	is a set of final states and $\Delta : Q \times \Sigma \rightarrow 2^Q$
	is a function of transitions.
\end{definition}

\begin{definition}
	A path $\pi$ in a finite-state automaton $\mathcal{A} = \langle
	\Sigma, Q, I, F, \Delta \rangle$ is a non-empty sequence $\langle q_1,
	\sigma_1, q_2 \rangle, \langle q_2, \sigma_2, q_3 \rangle, \dots,
	\langle q_{n-1}, \sigma_{n-1}, q_n \rangle$ where $\forall i \in
	\interval{1}{n}(q_i \in Q)$, $\forall i \in \interval{1}{n-1}(\sigma_i
	\in \Sigma)$ and $\forall i \in \interval{1}{n-1}(q_{i+1} \in
	\Delta(q_i, \sigma_i))$. The path $\pi$ is successful if $q_1 \in I$ and
	$q_n \in F$. The label of the path $\pi$ is the word
	$\sigma_1\sigma_2\ldots\sigma_{n-1}$ of length $n-1$ over $\Sigma$.
\end{definition}

\begin{definition}
	Let $\mathcal{A} = \langle \Sigma, Q, I, F, \Delta \rangle$ be a
	finite-state automaton. The generalized function of transitions
	$\Delta^* : Q \times \Sigma^* \rightarrow 2^Q$ is defined as
	the least function such that $\forall q_1 \in Q\; \forall \sigma \in
	\Sigma(q_2 \in \Delta(q_1, \sigma) \implies q_2 \in \Delta^*(q_1,
	\sigma))$ and $\forall q_1 \in Q\; \forall w \in \Sigma^*\; \forall
	\sigma \in \Sigma(q_2 \in \Delta^*(q_1, w) \land q_3 \in \Delta(q_2,
	\sigma) \implies q_3 \in \Delta^*(q_1, w\sigma))$.
\end{definition}

\begin{definition}
	The language recognised by a finite-state automaton $\mathcal{A} =
	\langle \Sigma, Q, I, F, \Delta \rangle$ is the set of all labels of
	successful paths $\mathcal{L}(\mathcal{A}) = \{ w \in \Sigma^* \mid
	\hat{\Delta}^*(I, \{ w \}) \cap F \neq \varnothing \}$.
\end{definition}

\begin{definition}
	A deterministic finite-state automaton $\mathcal{A}$ is a tuple $\langle
	\Sigma, Q, i, F,$\\
	$\delta \rangle$ where $\Sigma$ is an alphabet, $Q$ is a finite set of
	states,$i \in Q$ is an initial state, $F \subseteq Q$ is a set of final
	states and $\delta : Q \times \Sigma \rightarrow Q$ is a function of
	transitions.
\end{definition}

\begin{definition}
	A path $\pi$ in a deterministic finite-state automaton $\mathcal{A} =
	\langle \Sigma, Q, i, F, \delta \rangle$ is a non-empty sequence $\langle
	q_1, \sigma_1, q_2 \rangle, \langle q_2, \sigma_2, q_3 \rangle \dots,
	\langle q_{n-1}, \sigma_{n-1}, q_n \rangle$ where $\forall i \in
	\interval{1}{n}(q_i \in Q)$, $\forall i \in \interval{1}{n-1}(\sigma_i
	\in \Sigma)$ and $\forall i \in \interval{1}{n-1}(q_{i+1} = \delta(q_i,
	\sigma_i))$. The path $\pi$ is successful if $q_1 = i$ and $q_n \in F$.
	The label of the path $\pi$ is the word $\sigma_1\sigma_2\ldots\sigma_{n-1}$
	of length $n-1$ over $\Sigma$.
\end{definition}

\begin{definition}
	Let $\mathcal{A} = \langle \Sigma, Q, i, F, \delta \rangle$ be a
	deterministic finite-state automaton. The generalized function of
	transitions $\delta^* : Q \times \Sigma^* \rightarrow Q$ is defined as
	the least function such that $\forall q_1 \in Q\; \forall \sigma \in
	\Sigma(\delta(q_1, \sigma) = q_2 \implies \delta^*(q_1, \sigma) = q_2)$
	and $\forall q_1 \in Q\; \forall w \in \Sigma^*\; \forall \sigma \in
	\Sigma(\delta^*(q_1, w) = q_2 \land \delta(q_2, \sigma) = q_3 \implies
	\delta^*(q_1, w\sigma) = q_3)$.
\end{definition}

\begin{definition}
	The language recognised by a deterministic finite-state automaton
	$\mathcal{A} = \langle \Sigma, Q, i, F, \delta \rangle$ is the set of all
	labels of successful paths $\mathcal{L}(\mathcal{A}) = \{ w \in \Sigma^*
	\mid \delta^*(i, w) \in F \}$.
\end{definition}

\begin{theorem}[Rabin-Scott, see \cite{RS59}]
	Let $\mathcal{A} = \langle \Sigma, Q, I, F, \Delta \rangle$ be a
	finite-state automaton. We can construct a deterministic finite-state
	automaton $\mathcal{A'} = \langle Q, \Sigma', i, F', \delta \rangle$
	such that $\mathcal{L} ( \mathcal{A'} ) = \mathcal{L} ( \mathcal{A} )$
	where $Q' = 2^{Q}$, $F' = \{ q' \in Q' \mid q' \cap F \neq \varnothing
	\}$ and $\forall q' \in Q'\; \forall \sigma \in \Sigma(\delta(q',
	\sigma) = \bigcup_{q \in q'} \Delta^*(q, \sigma))$.
\end{theorem}

\begin{corollary}
	Let $\mathcal{A} = \langle \Sigma, Q, I, F, \Delta \rangle$ be a
	finite-state automaton and\\
	$\mathcal{A'} = \langle Q, \Sigma', i, F', \delta \rangle$ be the
	deterministic finite-state automaton derived from $\mathcal{A}$ by
	applying the Rabin-Scott construction. Then $| \delta | = 2^{|Q|} \times
	|\Sigma|$.
\end{corollary}

\begin{proof}
	$|\delta| = |Q'| \times |\Sigma| = |2^Q| \times |\Sigma| =
	2^{|Q|} \times |\Sigma|$.
\end{proof}

\begin{definition}
	The graph $\mathcal{G}_{\mathcal{A}}$ corresponding to a finite-state
	automaton $\mathcal{A} = \langle \Sigma, Q, I, F, \Delta \rangle$ is the
	pair $\langle Q, E \rangle$ where $E = \{ \langle q', q'' \rangle \mid
	q'' \in \Delta(q', \sigma) \}$.
\end{definition}

\begin{definition}
	A finite-state automaton $\mathcal{A}$ is planar if and only if its
	corresponding graph $\mathcal{G}_{\mathcal{A}}$ is planar.
\end{definition}
